\section{Related Work}

\subsection{Robot Learning}

Robot learning is a field of robotics that applies machine learning methods to robotic control tasks,
in hopes of producing systems that are more robust and general than classical hand engineered methods~\citep{capuano2025robot}.
Many recent methods have found success using end-to-end models, often referred to as policies, which directly map perception to actions using monolithic neural network architectures\cite{black2024pi0, ye2026dreamzero, jiang2026robotttcontextscaling, chi2025diffusion}.

These policies model a conditional distribution $\pi_\theta(a \mid o)$ of actions $a$ (eg. robot gripper pose, joint states, wheel velocity), given some observation $o$ (eg. camera input, joint states, depth maps, tactile inputs), such that the generated actions optimize some training objective. There are two prevalent approaches for training such robot learning policies: Reinforcement Learning and Behaviour Cloning, which primarily differ in their training objective and source of training data.

\subsubsection{Reinforcement Learning}

Reinforcement Learning (RL) methods work by training a policy to maximize some reward signal $r$ through trial-and-error interaction\cite{tang2025deep}.
This approach to robot learning allows for autonomous improvement of policies, as long as a reliable reward signal is available. This can prove difficult in the physical world, given that the world state needed for reward calculation needs to be inferred, unlike in simulated environments where world state is directly accessible. Furthermore, many RL algorithms require explicit exploration of their environments, which can be dangerous or time consuming on real hardware. Some other difficulties include reward hacking, reward specification, and sample inefficiency. Some of these obstacles have inspired work on sim2real methods, where models are first trained in simulation, and then adapted for use in the real world. %add more citations of RL projects 
It's important to note however, that a growing body of work has tested methods to get Reinforcement Learning working outside of hand coded simulators, whether that be in learned world models (section \ref{sec:worldmodels}), or in the physical world.

\subsubsection{Behaviour Cloning}

Many current Robot Learning systems utilize Behaviour Cloning (BC) to avoid the difficulties associated with RL. Behaviour Cloning methods are trained to minimize the difference between a policies action distribution and an expert distribution built from a dataset of expert actions $D$. These datasets are typically collected using direct teleoperation of a robot, however many recent works have made progress towards learning directly from human data and in-the-wild video data. %CITE
In parallel, large cross-embodiment datasets that pool demonstrations across many robots, tasks, and labs have made it possible to pre-train a single policy on far more data than any individual setup could produce, which has been a key enabler of generalist robot policies. %CITE Open X-Embodiment, AgiBot World, DROID, others
Although relatively simple to implement on real world tasks, Behaviour Cloning methods have their own distinct drawbacks. For example, models trained using BC can be susceptible to compounding error, causal confusion, mode averaging, over-fitting to expert demonstrations, and distributional shift. % CITE all
Despite these shortcomings, Behaviour Cloning methods, especially when paired with foundation models, have been the main driving force of recent advancements in robotic manipulation research.

\subsubsection{Visuo-motor policies and action distribution modelling}

Visuo-motor policies are a general category of models that take in, though not exclusively, visual input, and directly output robotic actions~\citep{chi2025diffusion}. A key design choice is how the action distribution $\pi_\theta(a \mid o)$ is modelled. The simplest approach is to regress actions directly with an MSE or L1 loss, but this assumes a unimodal distribution: when demonstrations contain several valid ways of completing a task, regression averages them into an action that may match none. Generative policies address this by learning the full distribution, either by discretizing actions into tokens and predicting them autoregressively %CITE RT-2, FAST
or by iteratively transforming noise into actions with diffusion~\citep{chi2025diffusion} or flow matching~\citep{lipman2023flowmatching}. Most modern policies also predict action chunks, sequences of $H$ future actions $a_{t:t+H}$, rather than a single step, which improves temporal consistency and reduces compounding error. %CITE ACT
% TBD A2A / much ado about noising

\subsubsection{Vision Language Action Models}
Vision-Language-Action models (VLAs) are a type of visuo-motor policy that additionally take in language inputs $l$, $\pi_\theta(a \mid o, l)$. Often, these models are built on top of pre-trained foundational Vision Language Models (VLMs), with the aim of transferring the semantic knowledge and instruction following ability learned from internet-scale image-text data to robotic control. The generic recipe takes a VLM backbone, feeds it camera images and a language instruction, and adds a mechanism for producing actions. Early VLAs represented actions as discrete tokens in the VLM's vocabulary, %CITE RT-2, OpenVLA
while more recent models such as $\pi_0$ attach a separate action expert that generates continuous action chunks with flow matching, conditioned on the VLM's internal representations~\citep{black2024pi0}. Follow-up work has improved open-world generalization by co-training on heterogeneous data sources %CITE pi0.5
and reduced the computational cost of these models to run on consumer hardware. %CITE SmolVLA
Despite this progress, VLAs have two key drawbacks. First, they are sample inefficient, often requiring hundreds of demonstrations to learn a single new task. %CITE
Second, they lack a physically grounded prior: their VLM backbones are trained primarily on static image-text pairs, which do not capture the dynamics needed for robotic control.


\subsubsection{World Action Models / Video Action Models}
World Action Models (WAMs, or sometimes called Video Action Models / VAMs) are a type of visuo-motor policy that predict both a chunk of future actions and the corresponding future observations, $\pi_\theta(a_{t:t+H}, o_{t+1:t+H} \mid o_t, l)$. Rather than starting from a VLM, these models are typically initialized from pre-trained latent video diffusion or flow matching transformers~\citep{peebles2023dit, lipman2023flowmatching}, and then fine-tuned on robot data to additionally produce actions~\citep{ye2026dreamzero}.

Existing WAMs can be broadly divided into two designs. In inverse dynamics (IDM) based approaches, a video model first predicts future observations, and a separate inverse dynamics model then infers the actions that would produce the predicted transition. %CITE UniPi, MimicVideo
This allows the video model to be trained on action-free video, with only the IDM requiring action-labelled robot data, but errors in the predicted video propagate directly into the actions. In joint approaches, a single model denoises video and action tokens together, so that action prediction attends directly to the predicted future~\citep{zhu2025uwm, ye2026dreamzero}.

The strong generalization reported for WAMs~\citep{ye2026dreamzero} has been attributed to two, not mutually exclusive, factors. The first is their prior: video generation models are pre-trained on web-scale video, and therefore already encode knowledge about object permanence, motion, and physical interaction. The second is the training objective: predicting future observations provides a dense supervisory signal that forces the model to learn scene dynamics, whereas action labels are low-dimensional and comparatively scarce. Disentangling the contribution of each remains an open question.
%CITE ablations on pre-training vs future prediction objective if available

The main drawback of WAMs is their computational cost, since generating future video requires many denoising steps over a high-dimensional latent, resulting in latency that is problematic for closed-loop control. %CITE FastWAM

\subsection{World Models}
\label{sec:worldmodels}
World Models (WMs) are action-conditioned models (not to be confused with action-generative WAMs) that predict how a world will evolve given some action sequence. This action sequence is not necessarily raw inputs to a robot, and can instead be some latent representation of meta-actions, language instructions, or any other modality that represents an action an agent takes to alter its environment. Similarly, the output of a World Model does not necessarily need to be visual, it can be any observed modality that represents world state (eg. tactile inputs, depth maps, audio, latent representations). 

A key property of World Models is that they are not action-generative like policies of WAMs. This means that in order to use a WM for control, it is required to search through various candidate action sequences, and compare the fitness, cost, or reward of each candidate. 

\subsubsection{Joint Embedding Predictive Architectures (JEPA)}
\label{sec:jepa}

Joint Embedding Predictive Architectures (JEPAs) are a self-supervised representation learning approach that makes predictions in latent space rather than in input space. A context encoder embeds an observed portion of the input, and a predictor is trained to infer the embedding of a target portion (e.g., masked image regions or future video frames~\citep{vjepa2_1}) as produced by the same encoder or an exponential moving average (EMA) copy of it. This differs from reconstruction-based methods, which must regenerate the target in its raw form (e.g., pixels or waveforms). Because the model only has to predict an abstract representation, it can discard unpredictable or irrelevant low-level detail and focus on more semantic features. To prevent the trivial solution where all inputs map to the same embedding, JEPAs typically use a stop-gradient on an EMA target encoder, or an explicit regularizer on the embedding distribution.
%CITE I-JEPA, V-JEPA2, LEJEPA

This approach to representation learning may be especially useful in robot learning, particularly in WAM settings, since much of the low-level detail that JEPAs can discard (e.g., texture, lighting, background) is irrelevant to control. However, most WAMs are initialized from pre-trained video diffusion models, which operate in the latent space of a VAE trained for pixel reconstruction, and therefore retain exactly this kind of detail. While JEPA-based action-conditioned world models have been explored~\citep{wang2026adajepa, maes2026lewm}, JEPA-based WAMs, which generate actions while predicting future latent states, remain largely unexplored.

\subsection{Continual Learning}


\subsubsection{Loss of Plasticity}

\subsubsection{TBD other continual learning literature}


